\section{The field emitted by a particle in an LG beam}

Consider a particle, initial in the position
\begin{align}
  {\bm r}(0) = {\boldsymbol{\mathfrak R}}_0
\end{align}
in the $Oxy$ plane. We denote by $(\rho_0, \varPhi_0)$ the cylindrical coordinates of the initial position of the particle.

The particle interacts with an LG beam of not too large intensity, of indices $(p,m)$. Then the trajectory pf the particle can be approximated according to the formulas in the previous section, with a $\xi$ dependent on the radial coordinate $\rho_0$.
\begin{importantbox}
  \begin{align}
    & {\bm r}_{0\perp}(\chi)=\frac{\xi (mc)}{\sqrt{2}k_L(n_L\cdot q_0)}\left[{\bm e}_x\sin(k_L\chi-m_L\varPhi_0)-{\bm e}_y\cos(k_L\chi-m_L\varPhi_0)\right]                                                  \qquad z_0(\chi)=\frac{q_0^3}{n\cdot q_0}\chi\\
    & {\boldsymbol{\beta}}_\perp(\phi)=\xi\frac{\sqrt{2}mc}{n_L\cdot q_0}\left[{\bm e}_x\cos(k_L\phi-m_L\varPhi_0)+{\bm e}_y\sin(k_L\phi-m_L\varPhi_0)\right]\qquad  \beta_z(\phi) =\frac{q_0^3}{q_0^0}
  \end{align}
\end{importantbox}

We use the previous results in the expression of the emitted electric field
\begin{align}
  {\bm E}({\bm x},k) =\frac{-i k e_0 e^{ik |{\bm x}_0|}}{4\pi\epsilon_0|{\bm x}_0|}\frac1{\sqrt{2\pi}}\int\limits_{-\infty}^{\infty}d(ct)\left[{{\bm n}_0\times({\bm n}_0\times{\boldsymbol\beta}(t))}\right]e^{ik(ct-{\bm n}_0\cdot{\bm r}_0(t))}
\end{align}
In the previous integral we perform another change of variables from $t$ to $\phi=ct-z$, using
\begin{align}
  d\phi = c (1-\beta_z(\phi)) dt = c\frac{n_L\cdot q_0}{q_0^0} dt,\qquad \rightarrow \qquad dt=\frac1c d\phi\frac{q_0^0}{n_L\cdot q_0}
\end{align}
\begin{warningbox}{}
  In the head-on geometry, low intensity limit, $\gamma\gg1$ the factor $\frac{q_0^0}{n_L\cdot q_0}\approx 2$
\end{warningbox}
After the change of variable form $t$ to $\phi$ we have
\begin{align}
  {\bm E}({\bm x},k) =\frac{-i k e_0 e^{ik |{\bm x}_0|}}{4\pi\epsilon_0|{\bm x}_0|}\frac1{\sqrt{2\pi}}\frac1{c}\frac{q_0^0}{n_L\cdot q_0}\int\limits_{-\infty}^{\infty}d\phi\left[{{\bm n}_0\times({\bm n}_0\times {\boldsymbol\beta}(\phi))}\right]e^{ik(\phi +z(\phi)-{{\bm n}_0\cdot{\bm r}_0(\phi)})}
\end{align}
and, denoting by $\theta_0$, $\phi_0$ the angles of ${\bm n}_0$ and ${\bm n}_L={\bm e}_z$
we have
\begin{align}
  z(\phi)-{\bm n}_0\cdot{\bm r}_0(\phi)={\bm n}_L\cdot{\bm r}_0(\phi)-{\bm n}_0\cdot{\bm r}_0(\phi)=({\bm n}_L-{\bm n}_0)\cdot{\bm r}_0(\phi)
\end{align}

Next we calculate $({\bm n}_L-{\bm n}_0)\cdot{\bm r}_0(\phi)$;
\begin{align}
  ({\bm n}_L-{\bm n}_0)\cdot{\bm r}_0(\phi)= & (1-\cos\theta_0)
  \frac{q_0^3}{n_L\cdot q_0} \phi     -\frac{\xi (mc)}{\sqrt{2}k_L(n_L\cdot q_0)}\sin\theta_0\left[\cos\phi_0\sin(k_L\phi-m_L\varPhi_0)-\sin\phi_0\cos(k_L\phi-m_L\varPhi_0)\right]\nonumber\\
  & (1-\cos\theta_0)
  \frac{q_0^3}{n_L\cdot q_0} \phi     -\frac{\xi (mc)}{\sqrt{2}k_L(n_L\cdot q_0)}\sin\theta_0\sin(k_L\phi-m_L\varPhi_0-\phi_0)
\end{align}
Using the four vectors
\begin{align}
  n_L=(1,{\bm n}_L)\equiv(1,{\bm e}_z),\qquad n_0=(1,{\bm n}_0)
\end{align}
we can write
\begin{align}
  1-\cos\theta_0=n_L\cdot n_0
\end{align}
Define
\begin{align}
  & \Gamma_0=\frac{\xi (mc)}{\sqrt{2}k_L(n_L\cdot)}\sin\theta_0,
\end{align}
\begin{align}
  \phi+({\bm n}_L-{\bm n}_0)\cdot{\bm r}_0(\phi)= & \phi\left(1+n_L\cdot n_0\frac{q_0^3}{n_L\cdot q_0}\right)-\Gamma_0\sin(k_L\phi-m_L\varPhi_0-\phi_0)
\end{align}
and using the Anger identity \cite{bessel}
\begin{align}
  e^{i z \sin \phi}=\sum_{n=-\infty}^{\infty} e^{i n \phi} J_n(z)
\end{align}
we have
\begin{align}
  e^{ik(\phi +z(\phi)-{{\bm n}_0\cdot{\bm r}_0(\phi)})} & =e^{ik\phi(1+n_L\cdot n_0\frac{q_0^3}{n_L\cdot q_0})}\sum\limits_{n=-\infty}^{\infty}J_n(k\Gamma_0)e^{-in(k_L\phi-m_L\varPhi_0-\phi_0)}
\end{align}
\begin{align}
  {\bm n}_0\times({\bm n}_0\times{\bm e}_{\{x,y,z\}})={\boldsymbol{\nu}}_{0\{x,y,z\}},\qquad {\boldsymbol{\nu}}_{\pm}=\zeta_x{\boldsymbol{\nu}}_{0x}\pm i\zeta_y{\boldsymbol{\nu}}_{0y},\qquad ({\boldsymbol{\nu}}_{\pm})^* = {\boldsymbol{\nu_{\mp}}}\label{defnupm}
\end{align}
and
\begin{align}
  {\bm n}_0\times({\bm n}_0\times{\boldsymbol{\beta}})  = & \xi\frac{mc}{n_L\cdot q_0}\frac1{2\sqrt{2}}\left[{\boldsymbol{\nu}}_{0x}\left(e^{i(k_L\phi -m_L\varPhi_0)}+e^{-i(k_L\phi-m_L\varPhi_0)}\right)-i{\boldsymbol{\nu}}_{0y}\left(e^{i(k_L\phi -m_L\varPhi_0)}-e^{-i(k_L\phi-m_L\varPhi_0)}\right)\right]+ \frac{q_0^3}{q_0^0}{\boldsymbol{\nu}_{0z}}                                                                                                                                              \\
  =                                                             & \xi\frac{mc}{n_L\cdot q_0}\frac1{2\sqrt{2}}{\boldsymbol{\nu}}_{-} e^{i(k_L\phi -m_L\varPhi_0)}+\xi\frac{mc}{n_L\cdot q_0}\frac1{2\sqrt{2}}{\boldsymbol{\nu}}_{+}e^{-(ik_L\phi -m_L\varPhi_0)}+{\boldsymbol{\nu}}_{0z}\frac{q_0^3}{q_0^0}
\end{align}

With these the Fourier transform of the electric field becomes
\begin{align}
  {\bm E}({\bm x},k) =\frac{-i k e_0 e^{ik |{\bm x}_0|}}{4\pi\epsilon_0|{\bm x}_0|}\frac1{\sqrt{2\pi}}\int\limits_{-\infty}^{\infty}d\phi & e^{ik\phi(1+n_L\cdot n_0\frac{q_0^3}{n_L\cdot q_0})}\sum\limits_{n=-\infty}^{\infty}J_n(k\Gamma_0)e^{in\phi_0}e^{-in(k_L\phi-m_L\varPhi_0)}\nonumber \\& \times\left\{\xi\frac{mc}{n_L\cdot q_0}\frac1{2\sqrt{2}}{\boldsymbol{\nu}}_{-} e^{i(k_L\phi -m_L\varPhi_0)}+\xi\frac{mc}{n_L\cdot q_0}\frac1{2\sqrt{2}}{\boldsymbol{\nu}}_{+}e^{-(ik_L\phi -m_L\varPhi_0)}+{\boldsymbol{\nu}}_{0z}\frac{q_0^3}{q_0^0}\right\}\nonumber\\
  & =\sum\limits_{K=-1,1}E_K({\bm x},k)
\end{align}
We also have
\begin{align}
  1+n_L\cdot n_0\frac{q_0^3}{n_L\cdot q_0} = \frac{n_L\cdot q_0 +(n_L\cdot n_0){q_0^3}}{n_L\cdot q_0} = \frac{q_0^0-q_0^3+(1-{\bf n}_L\cdot{\bf n}_0)q_0^3}{n_L\cdot q_0} = \frac{n_0\cdot q_0}{n_L\cdot q_0}
\end{align}
Then the terms in ${\bf E}_K({\bm x},k)$ can be written as
\begin{align}
  E_0({\bm x},k) = & \frac{q_0^3}{q_0^0}{\boldsymbol{\nu}_{0z}}   \frac{-i k e_0 e^{ik |{\bm x}_0|}}{4\pi\epsilon_0|{\bm x}_0|}\frac1{\sqrt{2\pi}}\int\limits_{-\infty}^{\infty}d\phi  e^{ik\phi\frac{n_0\cdot q_0}{n_L\cdot q_0}}\sum\limits_{n=-\infty}^{\infty}J_n(k\Gamma_0)e^{in\phi_0}e^{-in(k_L\phi-m_L\varPhi_0)}\nonumber                                                                 \\
  =& \frac{q_0^3}{q_0^0}{\boldsymbol{\nu}_{0z}}   \frac{-i k e_0 e^{ik |{\bm x}_0|}}{4\pi\epsilon_0|{\bm x}_0|}\frac1{\sqrt{2\pi}}\sum\limits_{n=-\infty}^{\infty}J_n(k\Gamma_0)e^{in\phi_0}e^{inm_L\varPhi_0}\int\limits_{-\infty}^{\infty}d\phi  e^{ik\phi\frac{n_0\cdot q_0}{n_L\cdot q_0}}e^{-ink_L\phi}\nonumber                                                                 \\
  & \frac{q_0^3}{q_0^0}{\boldsymbol{\nu}_{0z}}   \frac{-i k e_0 e^{ik |{\bm x}_0|}}{4\pi\epsilon_0|{\bm x}_0|}\frac1{\sqrt{2\pi}}\sum\limits_{n=-\infty}^{\infty}J_n(k\Gamma_0)e^{in\phi_0}e^{inm_L\varPhi_0}2\pi\delta(k\frac{n_0\cdot q_0}{n_L\cdot q_0}-nk_L)
\end{align}

\begin{align}
  E_{\pm1}({\bm x},k) = & \xi\frac{mc}{n_L\cdot q_0}\frac1{2\sqrt{2}}{\boldsymbol{\nu}_{\pm}}   \frac{-i k e_0 e^{ik |{\bm x}_0|}}{4\pi\epsilon_0|{\bm x}_0|}\frac1{\sqrt{2\pi}}\int\limits_{-\infty}^{\infty}d\phi  e^{ik\phi\frac{n_0\cdot q_0}{n_L\cdot q_0}}\sum\limits_{n=-\infty}^{\infty}e^{in\phi_0}J_n(k\Gamma_0)e^{-i(n\pm1)k_L\phi}e^{i(n\pm1)m_L\varPhi_0}\nonumber          \\
= & \xi\frac{mc}{n_L\cdot q_0}\frac1{2\sqrt{2}}{\boldsymbol{\nu}_{\pm}}   \frac{-i k e_0 e^{ik |{\bm x}_0|}}{4\pi\epsilon_0|{\bm x}_0|}\frac1{\sqrt{2\pi}}\sum\limits_{n=-\infty}^{\infty}e^{in\phi_0}J_n(k\Gamma_0)e^{i(n\pm1)m_L\varPhi_0}2\pi\delta(k\frac{n_0\cdot q_0}{n_L\cdot q_0}-(n\pm1)k_L))\nonumber          \\
\end{align}
With the notations
\begin{align}
k_N = N\frac{k_L}{\alpha},\qquad \alpha = 1+\frac{\xi^2}4n_0\cdot n_L
\end{align}We obtained
\begin{align}
\boxed{
  {\bm E}({\bm x},k)=\sum\limits_{N=-\infty}^{\infty}\tilde{\bm E}_N({\bm x})\delta(k-N \frac{k_L}{\alpha}),\label{Encirc}
}
\end{align}
with
\begin{align}
\boxed{
  \begin{aligned}
    \tilde{\bm E}_N({\bm x}) = & \frac{-i k_N  e_0 e^{ik_N (|{\bm x}_0|+Z_0)}}{4\pi\epsilon_0|{\bm x}_0|\alpha}{\sqrt{2\pi}}e^{iNm_L\varPhi_0}\left[{\boldsymbol{\nu}}_{+}\frac{\xi}{2}\tilde J_{N-1}(k_N;\gamma_0,\Gamma_0,\Delta_0)+{\boldsymbol{\nu}}_{-}\frac{\xi}{2}\tilde J_{N+1}(k_N;\gamma_0,\Gamma_0,\Delta_0)+\right.\nonumber \\
    & \left.{\boldsymbol{\nu}}_{0z}\left(\frac{\xi^2}4\tilde J_N(k_N;\gamma_0,\Gamma_0,\Delta_0)+\frac{\xi^2}8(\zeta_x^2-\zeta_y^2)\tilde J_{N-2}(k_N;\gamma_0,\Gamma_0,\Delta_0)+\frac{\xi^2}8(\zeta_x^2-\zeta_y^2)\tilde J_{N+2}(k_N;\gamma_0,\Gamma_0,\Delta_0)\right)\right],
  \end{aligned}
}
\end{align}

In the same approximation, the magnetic field  is
\begin{align}
\boxed{
{\bm B}({\bm x},k) = \frac1c{\bm n}_0\times{\bm E}({\bm x},k)=\sum\limits_{N=-\infty}^{\infty}\tilde {\bm B}_N({\bm x})\delta(k-N\frac{k_L}{\alpha})}\label{Bncirc}
\end{align}
with
\begin{equation}
\boxed{
  \begin{aligned}
    \tilde{\bm B}_N({\bm x}) = & \frac{-i k_N  e_0 e^{ik_N (|{\bm x}_0|+Z_0)}}{4\pi\epsilon_0|{\bm x}_0|c\alpha}{\sqrt{2\pi}}e^{iNm_L\varPhi_0}\left[{\boldsymbol{\mu}}_{+}\frac{\xi}{2}\tilde J_{N-1}(k_N;\gamma_0,\Gamma_0,\Delta_0)+{\boldsymbol{\mu}}_{-}\frac{\xi}{2}\tilde J_{N+1}(k_N;\gamma_0,\Gamma_0,\Delta_0)+\right.\nonumber \\
    & \left.{\boldsymbol{\mu}}_{0z}\left(\frac{\xi^2}4\tilde J_N(k_N;\gamma_0,\Gamma_0,\Delta_0)+\frac{\xi^2}8(\zeta_x^2-\zeta_y^2)\tilde J_{N-2}(k_N;\gamma_0,\Gamma_0,\Delta_0)+\frac{\xi^2}8(\zeta_x^2-\zeta_y^2)\tilde J_{N+2}(k_N;\gamma_0,\Gamma_0,\Delta_0)\right)\right],\nonumber
\end{aligned}}
\end{equation}

\begin{align}
\boxed{
  \begin{aligned}
    & {\bm e}_{\pm} = \zeta_x{\bm e}_x\pm i \zeta_y {\bm e}_y                                                                                                                                                                                                                                                                               \\
    & {\boldsymbol{\nu}}_{\pm} = {\bm n}_0\times\left({\bm n}_0\times{\bm e}_{\pm}\right)=-{\bm{n}}_0\times{\boldsymbol{\mu}_{\pm}},\qquad {                                                                       \boldsymbol{\nu}}_z={\bm n}_0\times\left({\bm n}_0\times{\bm e}_{z}\right)=-{\bm n}_0\times{\boldsymbol{\mu}}_z\nonumber \\
    & {\boldsymbol{\mu}}_{\pm}={\bm n}_0\times{\boldsymbol{\nu}}_{\pm}=-{\bm n}_0\times{\bm e}_{\pm},\qquad {\boldsymbol{\mu}}_{0z}={\bm n}_0\times{\boldsymbol{\nu}}_{0z}=-{\bm n}_0\times{\bm e}_z                                                                                                                 \nonumber              \\
    & \tilde J_M(k;\gamma_0, \Gamma_0, \Delta_0) = \sum\limits_{p=-\infty}^{\infty}e^{i(M+2p)\gamma_0}J_{M+2p}(k\Gamma_0)J_p(k\Delta_0) = \sum\limits_{p=-\infty}^{\infty}e^{i(M-2p)\gamma_0}J_{M-2p}(k\Gamma_0)J_p(k\Delta_0) (-1)^p,\nonumber                                                                                             \\
    & \gamma_0=\arctan\left(\frac{\sin\phi_0\zeta_y}{\cos\phi_0\zeta_x}\right),\qquad \Gamma_0=\frac{\xi}{k_L}\sin\theta_0\sqrt{\cos^2\phi_0\zeta_x^2+\sin^2\phi_0\zeta_y^2},\qquad \Delta_0 = n_L\cdot n_0\frac{\xi^2}{8k_L}(\zeta_x^2-\zeta_y^2)\nonumber                                                                                 \\
    & (1-\cos\theta_0)z_i = n_L\cdot n_0 \frac{\xi^2}4\frac{\zeta_x^2-\zeta_y^2}{2k_L}\sin(2m_L\varPhi_0)=Z_0\nonumber
\end{aligned}}
\end{align}
The conservation law is consistent with the general formula of non-linear Thomson scattering
\begin{align}
N k_L\cdot q = k\cdot  q,\qquad q = p + mc\frac{\xi^2}4 n_L
\end{align}
which for the case of an electron initially at rest becomes
\begin{align}
N k_L mc = k mc (1+\frac{\xi^2}4n_0\cdot n_L),\qquad k = \frac{Nk_L}{1+\frac{\xi^2}4n_0\cdot n_L}
\end{align}
The previous result shows that if the retardation is included the emitted frequency depends on the direction; however, for the geometry considered in our case the dependence is very weak.

Properties of the generalized Bessel functions
\begin{align}
\tilde J_M(k;\gamma_0, \Gamma_0, \Delta_0)   & = \sum\limits_{p=-\infty}^{\infty}e^{i(M+2p)\gamma_0}J_{M+2p}(k\Gamma_0)J_p(k\Delta_0)                                                                                                       \\
\tilde J_{-M}(-k;\gamma_0,\Gamma_0,\Delta_0) & = \sum\limits_{p=-\infty}^{\infty}e^{i(-M+2p)\gamma_0}J_{-M+2p}(-k\Gamma_0)J_p(-k\Delta_0)=\sum\limits_{p=-\infty}^{\infty}e^{i(-M+2p)\gamma_0}J_{M-2p}(k\Gamma_0)J_{-p}(k\Delta_0)\nonumber \\
& =\sum\limits_{p=-\infty}^{\infty}e^{-i(M+2p)\gamma_0}J_{M+2p}(k\Gamma_0)J_{p}(k\Delta_0)=[\tilde J_{M}(k;\gamma_0,\Gamma_0,\Delta_0)]^*
\end{align}
Then we can easily prove that
\begin{align}
{\bm E}_{-N}({\bm x}) = (E_N({\bm x}))^*,\qquad {\bm B}_{-N}({\bm x}) = (B_N({\bm x}))^*
\end{align}

\section{Approximated formulas}

In our calculation $\theta_0$ is of the order of $\frac{\lambda}{w_0}$, we consider only the terms of the lowest order in $\theta_0$.
\begin{align}
\Gamma_0 = \frac{\xi}{k_L} \theta_0\sqrt{cos^2\phi_0\zeta_x^2+\sin^2\phi_0\zeta_y^2},\qquad \Delta_0=\frac{\xi^2}{8k_L}(\zeta_x^2-\zeta_y^2)(1-\cos\theta_0 ) = \frac{\xi^2}{16k_L}\theta_0^2(\zeta_x^2-\zeta_y^2)
\end{align}

\subsubsection{Approximation of the generalized Bessel function}

\begin{align}
& \tilde J_M(k;\gamma_0, \Gamma_0, \Delta_0) = \sum\limits_{p=-\infty}^{\infty}e^{i(M+2p)\gamma_0}J_{M+2p}(k\Gamma_0)J_p(k\Delta_0) = \sum\limits_{p=-\infty}^{\infty}e^{i(M-2p)\gamma_0}J_{M-2p}(k\Gamma_0)J_p(k\Delta_0) (-1)^p,\nonumber \\
& J_n(x)\approx \frac{x^n}{2^n n!}\quad {\text{for   }}n\ge0 ,\qquad J_{-n}(x)=(-1)^n J_n(x)
\end{align}
\begin{enumerate}

\item the case $M=-1$
  \begin{align}
    \tilde J_{-1}(k;\gamma_0, \Gamma_0, \Delta_0) = \sum\limits_{p=-\infty}^{\infty}e^{i(-1+2p)\gamma_0}J_{-1+2p}(k\Gamma_0)J_p(k\Delta_0)
  \end{align}
  the lowest order term is the term with $p=0$
  \begin{align}
    \tilde J_{-1}(k;\gamma_0, \Gamma_0, \Delta_0) = e^{i(-1)\gamma_0}J_{-1}(k\Gamma_0)=- \frac{e^{-i\gamma_0}}2(k\Gamma_0)
  \end{align}
\item the case $M=0$
  \begin{align}
    \tilde J_0(k;\gamma_0, \Gamma_0, \Delta_0) = \sum\limits_{p=-\infty}^{\infty}e^{2ip\gamma_0}J_{2p}(k\Gamma_0)J_p(k\Delta_0)\approx 1
  \end{align}
\item the case $M=1$
  \begin{align}
    \tilde J_1(k;\gamma_0, \Gamma_0, \Delta_0) = \sum\limits_{p=-\infty}^{\infty}e^{i(1-2p)\gamma_0}J_{1-2p}(k\Gamma_0)J_p(k\Delta_0) (-1)^p
  \end{align}
  The lowest order term is the term with $p=0$
  \begin{align}
    \tilde J_1(k;\gamma_0, \Gamma_0, \Delta_0) \approx e^{i\gamma_0}J_{1}(k\Gamma_0)\approx  \frac{e^{i\gamma_0}}{2}(k\Gamma_0)
  \end{align}
\item the case $M\ge2$
  \begin{align}
    \tilde J_M(k;\gamma_0, \Gamma_0, \Delta_0) = \sum\limits_{p=-\infty}^{\infty}e^{i(M-2p)\gamma_0}J_{M-2p}(k\Gamma_0)J_p(k\Delta_0) (-1)^p
  \end{align}
  To the previous sum contribute the terms with $p\le M/2$
  \begin{align}
    \tilde J_M(k;\gamma_0, \Gamma_0, \Delta_0) \approx & \sum\limits_{p=0}^{M/2}e^{i(M-2p)\gamma_0}\frac{(k\Gamma_0)^{M-2p}}{2^{M-2p}(M-2p)!}\frac{(k\Delta_0)^p}{2^p p!}(-1)^p\nonumber                                                           \\
    =                                                  & \xi^M\theta_0^M\sum\limits_{p=0}^{M/2}\frac{e^{i(M-2p)\gamma_0}(-1)^p(k/k_L)^{M-p}}{2^{M+3p}(M-2p)!p!}({\cos^2\phi_0\zeta_x^2+\sin^2\phi_0\zeta_y^2})^{(M-2p)/2}(\zeta_x^2-\zeta_y^2)^{p}
  \end{align}

\end{enumerate}

\subsubsection{Approximation ofthe unity vectors}

We also approximate the unity vectors $\mu_{\pm}$, $\nu_{\pm}$, $\mu_z$, $\nu_z$. Up to the lowest order in $\theta_0$ we have
\begin{align}
& {\boldsymbol \mu}_+=[i\zeta_y,-\zeta_x,0]+{\cal O}(\theta_0)                           \\
& {\boldsymbol \mu}_-=[-i\zeta_y,-\zeta_x,0]+{\cal O}(\theta_0)                          \\
& {\boldsymbol \mu}_{0z}=[-\theta_0\sin\phi_0,\theta_0\cos\phi_0,0]+{\cal O}(\theta_0^2) \\
& {\boldsymbol \nu}_+=[-\zeta_x,-i\zeta_y,0]+{\cal O}(\theta_0)                          \\
& {\boldsymbol \nu}_-=[-\zeta_x,i\zeta_y,0]+{\cal O}(\theta_0)                           \\
& {\boldsymbol \nu}_{0z}=[\theta_0\cos\phi_0,\theta_0\sin\phi_0,0]+{\cal O}(\theta_0^2)
\end{align}

Within this approximation, up to the lowest order in  $\theta_0$, the emitted field becomes for $N\ge2$

\begin{empheq}[box=\fbox]{align}
\tilde{\bm E}_N({\bm x}) = & \frac{-i k_N  e_0 e^{ik_N (|{\bm x}_0|+Z_0)}}{4\pi\epsilon_0|{\bm x}_0|}{\sqrt{2\pi}}e^{iNm_L\varPhi_0}\left[{\boldsymbol{\nu}}_{+}\frac{\xi}{2}\tilde J_{N-1}(k_N;\gamma_0,\Gamma_0,\Delta_0)+{\boldsymbol{\nu}}_{0z}\frac{\xi^2}8(\zeta_x^2-\zeta_y^2)\tilde J_{N-2}(k_N;\gamma_0,\Gamma_0,\Delta_0)\right],\nonumber \\&=\frac{-i k_N  e_0 e^{ik_N (|{\bm x}_0|+Z_0)}}{4\pi\epsilon_0|{\bm x}_0|}{\sqrt{2\pi}}e^{iNm_L\varPhi_0}\nonumber\\
& \qquad\;\left[{\bm e}_x\left(-\zeta_x\frac{\xi}2J_{N-1}(k_N;\gamma_0,\Gamma_0,\Delta_0)+\theta_0\cos\phi_0\frac{\xi^2}8(\zeta_x^2-\zeta_y^2)\tilde J_{N-2}(k_N;\gamma_0,\Gamma_0,\Delta_0)\right)\right.\nonumber                                                                                                        \\
& \qquad +\left.{\bm e}_y\left(-i\zeta_y\frac{\xi}2J_{N-1}(k_N;\gamma_0,\Gamma_0,\Delta_0)+\theta_0\sin\phi_0\frac{\xi^2}8(\zeta_x^2-\zeta_y^2)\tilde J_{N-2}(k_N;\gamma_0,\Gamma_0,\Delta_0)\right)\right]
\end{empheq}
\begin{empheq}[box=\fbox]{align}
\tilde{\bm B}_N({\bm x}) = & \frac{-i k_N  e_0 e^{ik_N (|{\bm x}_0|+Z_0)}}{4\pi\epsilon_0|{\bm x}_0| c}{\sqrt{2\pi}}e^{iNm_L\varPhi_0}\left[{\boldsymbol{\mu}}_{+}\frac{\xi}{2}\tilde J_{N-1}(k_N;\gamma_0,\Gamma_0,\Delta_0)+{\boldsymbol{\mu}}_{0z}\frac{\xi^2}8(\zeta_x^2-\zeta_y^2)\tilde J_{N-2}(k_N;\gamma_0,\Gamma_0,\Delta_0)\right],\nonumber \\&=\frac{-i k_N  e_0 e^{ik_N (|{\bm x}_0|+Z_0)}}{4\pi\epsilon_0|{\bm x}_0|}{\sqrt{2\pi}}e^{iNm_L\varPhi_0}\nonumber\\
& \qquad\;\left[{\bm e}_x\left(i\zeta_y\frac{\xi}2J_{N-1}(k_N;\gamma_0,\Gamma_0,\Delta_0)-\theta_0\sin\phi_0\frac{\xi^2}8(\zeta_x^2-\zeta_y^2)\tilde J_{N-2}(k_N;\gamma_0,\Gamma_0,\Delta_0)\right)\right.\nonumber                                                                                                          \\
& \qquad +\left.{\bm e}_y\left(-\zeta_x\frac{\xi}2J_{N-1}(k_N;\gamma_0,\Gamma_0,\Delta_0)+\theta_0\cos\phi_0\frac{\xi^2}8(\zeta_x^2-\zeta_y^2)\tilde J_{N-2}(k_N;\gamma_0,\Gamma_0,\Delta_0)\right)\right]\nonumber                                                                                                          \\
& =\frac1c{\bm e}_z\times \tilde{\bm E}_N({\bm x})
\end{empheq}

and for $N=1$
\begin{empheq}[box=\fbox]{align}
\tilde{\bm E}_1({\bm x}) = & \frac{-i k_1  e_0 e^{ik_1 (|{\bm x}_0|+Z_0)}}{4\pi\epsilon_0|{\bm x}_0|}{\sqrt{2\pi}}e^{im_L\varPhi_0}\left[{\boldsymbol{\nu}}_{+}\frac{\xi}{2}\tilde J_{0}(k_1;\gamma_0,\Gamma_0,\Delta_0)\right],\nonumber \\&=\frac{-i k_N  e_0 e^{ik_N (|{\bm x}_0|+Z_0)}}{4\pi\epsilon_0|{\bm x}_0|}{\sqrt{2\pi}}e^{im_L\varPhi_0}\left[{\bm e}_x\left(-\zeta_x\frac{\xi}2J_{0}(k_N;\gamma_0,\Gamma_0,\Delta_0)\right)+{\bm e}_y\left(-i\zeta_y\frac{\xi}2J_{0}(k_1;\gamma_0,\Gamma_0,\Delta_0)\right)\right]
\end{empheq}
\begin{empheq}[box=\fbox]{align}
\tilde{\bm B}_1({\bm x}) = & \frac{-i k_1  e_0 e^{ik_1 (|{\bm x}_0|+Z_0)}}{4\pi\epsilon_0|{\bm x}_0|c}{\sqrt{2\pi}}e^{im_L\varPhi_0}\left[{\boldsymbol{\mu}}_{+}\frac{\xi}{2}\tilde J_{0}(k_1;\gamma_0,\Gamma_0,\Delta_0)\right],\nonumber \\&=\frac{-i k_N  e_0 e^{ik_N (|{\bm x}_0|+Z_0)}}{4\pi\epsilon_0|{\bm x}_0|}{\sqrt{2\pi}}e^{im_L\varPhi_0}\left[{\bm e}_x\left(i\zeta_y\frac{\xi}2J_{0}(k_N;\gamma_0,\Gamma_0,\Delta_0)\right)+{\bm e}_y\left(\zeta_x\frac{\xi}2J_{0}(k_1;\gamma_0,\Gamma_0,\Delta_0)\right)\right]\nonumber                                                                                                          \\
& =\frac1c{\bm e}_z\times \tilde{\bm E}_1({\bm x})
\end{empheq}
where
\begin{align}
k_N = {Nk_L}({1-\frac{\xi^2}{8}\theta_0^2})\approx{Nk_L}
\end{align}

\subsection{Circular Polarization}
In the circular polarization case the previous formulas reduce to
\begin{align}
\tilde J_M(k;\gamma_0, \Gamma_0, \Delta_0) \approx  \xi^M\theta_0^M\frac{e^{iM\gamma_0}(k/k_L)^{M}}{2^{3M/2}M!},\qquad Z_0=0
\end{align}
\begin{align}
\tilde{\bm E}_N({\bm x}) \approx & \frac{-i k_N  e_0 e^{ik_N (|{\bm x}_0|)}}{4\pi\epsilon_0|{\bm x}_0|}{\sqrt{2\pi}}e^{iNm_L\varPhi_0}\left(-\zeta_x{\bm e}_x-i\zeta_y{\bm e}_y\right)\frac{\xi}2\tilde J_{N-1}(k_N;\gamma_0,\Gamma_0,\Delta_0)\nonumber                    \\
\approx                          & \frac{i k_N  e_0 e^{ik_N |{\bm x}_0|}}{4\pi\epsilon_0|{\bm x}_0|}{\sqrt{2\pi}}e^{iNm_L\varPhi_0}\left(\zeta_x{\bm e}_x+i\zeta_y{\bm e}_y\right)\frac{\xi^N}2 \theta_0^{N-1}\frac{e^{i(N-1)\gamma_0}N^{N-1}}{2^{3(N-1)/2}(N-1)!}\nonumber \\
=                                & \frac{i k_L  e_0}{4\pi\epsilon_0}\frac{N^N}{(N-1)!2^{3N/2-1/2}}\sqrt{2\pi}(\zeta_x{\bm e}_x+i \zeta_y{\bm e}_y)\frac{e^{iNk_L|{\bm x}_0|}}{|{\bm x}_0|}\xi^N\theta_0^{N-1}e^{i(Nm_L\varPhi_0+(N-1)\gamma_0)}\label{approx-e-n-circ}
\end{align}
\begin{align}
\tilde{\bm B}_N({\bm x})= \frac1c{\bm e}_z\times\tilde{\bm E}_N({\bm x})
\end{align}

\section{Integration over the initial electron distribution}

We consider the case of a charge distributed uniformly with the charge density $\sigma$ within a disk of radius $\rho_m$ in the plane $Oxy$, interacting with an $LG$ beam. We assume that the beam intensity is low, such that the amplitude of the oscillatory motion of the particle is low with respect to waist $w_0$ of the external beam. we also assume that the particles are initially at rest, in the focal plane $Oxy$.

A particle with the initial position $\rho_0, \varPhi_0$ will interact with the field

\begin{align}
& E_x({\bm r},t) = E_0\Re\left\{\zeta_x u_{p_Lm_L}(\rho_0)e^{-i(\omega t-k_Lz-m_L\varPhi_0}\right\},\qquad
  E_y({\bm r},t) = E_0\Re\left\{\zeta_y u_{p_Lm_L}(\rho_0)e^{-i(\omega t-k_Lz-m_L\varPhi_0)}\right\}          \\
  & {\bm B}({\bm r},t) = {\bm e}_z\times{\bm E}({\bm r},t)
\end{align}
with
\begin{align}
  u_{pm}(\rho_0,\varPhi_0) = N_{pm}\left(\frac{\sqrt{2}}{w_0}\right)^{|m|}\rho_0^{|m|}e^{-\frac{x^2+y^2}{w_0^2}}\, _1F_1(-p, |m|+1, \frac{2\rho_0^2}{w_0^2})
  \qquad      N_{pm} = \frac{\sqrt{2}}{|m|!}\sqrt{\frac{(p+|m|)!}{p!}}
\end{align}

Then the field emitted by this distribution will be calculated as the superposition of the fields emitted by the charge in each surface element within the disk.

\begin{align}
  {\bm E}_{tot}({\bm x},t) = \int \limits_{\Sigma} da {\bm E}({\bm x},t)
\end{align}
and the Fourier transform becomes
\begin{align}
  {\bm E}_{tot}({\bm x},\omega) = \sum\limits_{N=-\infty}^{\infty} \tilde{\bm E}_{N,tot}({\bm x})\delta(\omega-N\omega_L)
\end{align}
where
\begin{align}
  \tilde {\bm E}_{N,tot}({\bm x})=\int\limits_0^{\rho_m}d\rho_0\rho_0\int\limits_0^{2\pi}d\varPhi_0 \tilde{\bm E}_N({\bm x})\label{int-En}
\end{align}
Similarly, the Fourier transform of the total magnetic field radiated is
\begin{align}
  {\bm B}_{tot}({\bm x},\omega) = \sum\limits_{N=-\infty}^{\infty} \tilde{\bm B}_{N,tot}({\bm x})\delta(\omega-N\omega_L) ,\qquad    \tilde {\bm B}_{N,tot}({\bm x})=\int\limits_0^{\rho_m}d\rho_0\rho_0\int\limits_0^{2\pi}d\varPhi_0 \tilde{\bm B}_N({\bm x})
\end{align}
In the equations above $\tilde{\bm E}_N$, $\tilde {\bm B}_N$ are in (\ref{Encirc},\ref{Bncirc}) above.

The  energy passing through a surface $\Sigma$ is
\begin{align}
  W_{\Sigma} & =\int\limits_{-\infty}^\infty dt \int\limits_\Sigma d{\bm a}\cdot{\bm S}({\bm x},t)\nonumber                                                                                                                                                                \\
  & = \int\limits_\Sigma d{\bm a}\cdot\int\limits_{-\infty}^\infty dt\frac1{\mu_0}{\bm E({\bm x},t)}\times{\bm B}^*({\bm x},t)\nonumber                                                                                                                         \\
  & = \int\limits_\Sigma d{\bm a}\cdot\int\limits_{-\infty}^\infty dt\frac1{\mu_0}\frac1{2\pi}\int \limits_{-\infty}^\infty d\omega \int \limits_{-\infty}^\infty d\omega' e^{i\omega t} e^{-i\omega' t}{\bm E({\bm x},\omega)}\times{\bm B}^*({\bm x},\omega')
\end{align}
In the previous equations we have used that fact that the fields are real, so ${\bm B}^*={\bm B}$. Next we change the order of integration and use the $\delta$ function representation
\begin{align}
  \delta(\omega-\omega')=\frac1{2\pi}\int\limits_{-\infty}^{\infty} dt e^{i(\omega-\omega')t}
\end{align}
with the result
\begin{align}
  W_{\Sigma} & =\int\limits_{\Sigma} d{\bm a}\int\limits_{-\infty}^{\infty} d\omega \frac1{\mu_0}E({\bm x},\omega)\times{\bm B}({\bm x},\omega)=\int\limits_{\Sigma} d{\bm a}\int\limits_{-\infty}^{\infty} d\omega \sum\limits_{N}\sum\limits_{N'}\frac1{\mu_0}{\bm E}_N({\bm x})\times{\bm B}^*_{N'}({\bm x})\delta(\omega-N\omega L)\delta(\omega-N'\omega_L)\nonumber \\
  & =\sum\limits_N W^{(N)}_{\Sigma} T_t
\end{align}
with
\begin{align}
  W^{(N)}_\Sigma = \int\limits_{\Sigma} d{\bm a}\frac1{\mu_0} {\bm E}_N({\bm x})\times{\bm B}^*_N({\bm x})
\end{align}
and $T_t$ the total duration of interaction.

The $W^{(N)}_\Sigma$ is the flux of the $\omega_N$ component of the radiated field.

\subsection{Analytical estimation of the phase of the emitted field in the circular polarization case}

We consider the expression of the total emitted electric field of given harmonic index $N$ (\ref{int-En})
\begin{align}
  \tilde {\bm E}_{N,tot}({\bm x})=\int\limits_0^{\rho_m}d\rho_0\rho_0\int\limits_0^{2\pi}d\varPhi_0 \tilde{\bm E}_N({\bm x})
\end{align}
and use here the approximation of $\tilde{\bm E}_N$ (\ref{approx-e-n-circ})
\begin{align}
  \tilde {\bm E}_{N,tot}({\bm x}) & =\int\limits_0^{\rho_m}d\rho_0\rho_0\int\limits_0^{2\pi}d\varPhi_0  \frac{i k_L  e_0}{4\pi\epsilon_0}\frac{N^N}{(N-1)!2^{3N/2-1/2}}\sqrt{2\pi}(\zeta_x{\bm e}_x+i \zeta_y{\bm e}_y)\frac{e^{iNk_L|{\bm x}_0|}}{|{\bm x}_0|}\xi(\rho_0)^N\theta_0^{N-1}e^{i(Nm_L\varPhi_0+(N-1)\gamma_0)}\nonumber \\
  & =\frac{i k_L  e_0}{4\pi\epsilon_0}\frac{N^N}{(N-1)!2^{3N/2-1/2}}\sqrt{2\pi}(\zeta_x{\bm e}_x+i \zeta_y{\bm e}_y)\int\limits_0^{\rho_m}d\rho_0\rho_0\xi^N(\rho_0)\int\limits_0^{2\pi}d\varPhi_0  \frac{e^{iNk_L|{\bm x}_0|}}{|{\bm x}_0|}\theta_0^{N-1}e^{i(Nm_L\varPhi_0+(N-1)\gamma_0)}\nonumber \\
\end{align}
Next we analyze  the integral over $\varPhi_0$; let us denote the coordinates of the observation point by
\begin{align}
  {\bm x}=R\cos\phi{\bm e}_x+R\cos\phi{\bm e}_y+Z{\bm e}_z
\end{align}
In a typical calculation $Z$ is fixed and $R$, $\phi$ changes (we record the radiation on a screen perpendicular on the beam propagation direction at the distance $Z$ from the focal plane).
The variable $|x_0|$  becomes
\begin{align}
  |{\bm x}_0|=\sqrt{(R\cos\phi-\rho_0\cos\varPhi_0)^2+(R\cos\phi-\rho_0\cos\varPhi_0)^2+Z^2}
\end{align}
as $e^{iNk_L|{\bm x}_0|}$ is a very fast oscillating function we can use the stationary phase method for estimating the integral.
We use the general formula \cite{ref2}
\begin{align}
  {\cal I}_{\varPhi_0}= \int d\varPhi_0 e^{if(\varPhi_0)} g(\varPhi_0)\approx\sum\limits_{x_0}\sqrt\frac{2\pi}{|f''(x_0)|}g(x_0)e^{if(x_0)+i\sigma\frac{\pi}4}
\end{align}
where $x_0$ {are} the stationary points of $f(\varPhi_0)$
\begin{align}
  f'(x_0)=0
\end{align}
and $\sigma$ is the sign of the second derivative $f''(x_0)$

In our case the derivative of  $f(\varPhi_0)$ is
\begin{align}
  f'(\varPhi_0)=-\frac{k R \rho \sin(\phi-\varPhi_0)}{\sqrt{R^2+Z^2+\rho^2-2 R \rho \cos(\phi-\varPhi_0)}}
\end{align}
whose zeroes are
\begin{align}
  \Phi_0^{(1)} = \phi,\qquad \Phi^{(2)}_0=\phi+\pi
\end{align}
and the integral over $\varPhi_0$ becomes
\begin{align}
  & {\cal I}_{\varPhi_0}=e^{i(m_LN\phi+\epsilon(N-1)\phi)}h(\rho_0),\nonumber                                                                                                              \\
  & h(\rho_0)=\sqrt{2\pi i}\left[e^{iNk_L\sqrt{Z^2+(R-\rho_0)^2}}\left(\frac{R-\rho_0}{\sqrt{Z^2+(R-\rho_0)^2}}\right)^N\sqrt{\frac{\sqrt{Z^2+(R-\rho_0)^2}}{Nk_LR\rho_0}}\right.\nonumber \\
  & \qquad\qquad\left.-i(-1)^{Nm_L}e^{iNk_L\sqrt{Z^2+(R+\rho_0)^2}}\left(\frac{R+\rho_0}{\sqrt{Z^2+(R+\rho_0)^2}}\right)^N\sqrt{\frac{\sqrt{Z^2+(R+\rho_0)^2}}{Nk_LR\rho_0}}\right]
\end{align}
and $\epsilon$ is the helicity of the circular polarization $\epsilon={\mathrm sgn}(\zeta_x/\zeta_y)$

With these the expression of the total  electric field measured in ${\bm x}$ is
\begin{align}
  \tilde {\bm E}_{N,tot}({\bm x}) & =\frac{i k_L  e_0}{4\pi\epsilon_0}\frac{N^N}{(N-1)!2^{3N/2-1/2}}\sqrt{\pi}({\bm e}_x+i \epsilon{\bm e}_y)\int\limits_0^{\rho_m}d\rho_0\rho_0\xi^N(\rho_0)\int\limits_0^{2\pi}d\varPhi_0  \frac{e^{iNk_L|{\bm x}_0|}}{|{\bm x}_0|}\theta_0^{N-1}e^{i(Nm_L\varPhi_0+(N-1)\gamma_0)}\nonumber
\end{align}
